% cover_letter.tex — Cover letter for IEEE Access
% Compile: bash build.sh cover
\documentclass[11pt,a4paper]{article}
\usepackage[margin=2.5cm]{geometry}
\usepackage[utf8]{inputenc}
\usepackage[T1]{fontenc}
\usepackage{lmodern}
\usepackage{amsmath,amssymb}
\usepackage{hyperref}
\hypersetup{hidelinks}

\setlength{\parindent}{0pt}
\setlength{\parskip}{0.6em}

\begin{document}

\begin{center}
{\Large\bfseries Cover Letter}\\[4pt]
{\itshape Original submission to \emph{IEEE Access}}
\end{center}

\vspace{1em}

\noindent\textbf{Date:} \today\\[2pt]
\textbf{To:} The Editor-in-Chief, \emph{IEEE Access}\\
\textbf{From:} Bin Nong, Weihong Fu, and Zhuoyun Jiang\\
\textbf{Corresponding email:} \href{mailto:nongbin@tfswufe.edu.cn}{nongbin@tfswufe.edu.cn}\\
\textbf{Manuscript title:} On the Waveform-to-Bit Gap in Single-Channel
Blind Source Separation of Co-Frequency Signals: A Decision-Identical
Evaluation Protocol and Controlled Diagnosis\\
\textbf{Manuscript type:} Original Research Article

\vspace{1em}

\noindent Dear Editor,

\vspace{0.4em}

We are pleased to submit the above manuscript for consideration in
\emph{IEEE Access}. The paper contributes \emph{measurement
infrastructure and a controlled diagnosis} for a question that the
single-channel blind source separation (SC-BSS) literature has so far
left untested: whether the waveform-fidelity gains that the field
optimises and reports (SI-SDR/SDR) actually reach the demodulated bits
of co-frequency communication signals. The work sits squarely within
the journal's scope in signal processing, communications, and machine
learning for the physical layer, and it is designed for reuse: the
benchmark, the evaluation receiver, and the training-time proxy are
released in full.

\paragraph{Key contributions.}
\begin{enumerate}
\item \textbf{A decision-identical evaluation protocol.} All bit-level
claims in the literature are only as strong as the receiver that
produces them. We contribute (i) an offset-compensated receiver
(oracle down-conversion, matched filtering, symbol-grid sampling,
unit-power normalisation, phase-only alignment) producing compensated
SER and Gray-mapped BER from shared decisions, and (ii) a
differentiable soft-demodulation loss whose front-end is verified
\emph{symbol-for-symbol identical} to that evaluation receiver
(matched filter agreeing with the reference implementation to a
relative error of $3.5\times10^{-7}$). This removes the train/eval
mismatch that has made earlier SER claims hard to compare, and it is
directly reusable wherever separation is evaluated against downstream
bits.
\item \textbf{A quantified waveform-to-bit mismatch with a located
cause.} On a variable-count ($K \in \{1,2,3\}$) co-frequency benchmark
at $\mathrm{SIR} \approx 0$\,dB, a $+4.1$\,dB SI-SDR improvement moves
compensated SER only from 0.5049 to 0.5027---even $+8.4$\,dB at
$-10$\,dB SNR buys nothing---and the flatness persists across an SIR
sweep to $20$\,dB at fixed $K{=}2$. Per-$K$ decomposition localises
the damage: a single interferer raises the baseline SER $3.8\times$
(0.146 to 0.561), so the gap is interference-limited, not
noise-limited. To our knowledge this mismatch has not previously been
quantified under statistically controlled conditions.
\item \textbf{A sign-structured effect of the task-oriented loss.} The
soft-demodulation term is practically negligible in the pool ($\leq
0.03$\,pp in the fair additive comparison), but the pool hides a
reproducible structure: the term \emph{helps} at $K{=}2$ (5/5 seeds
positive in all four contrasts) and \emph{hurts} at $K{=}3$ (0/5
seeds)---the number of interferers, not the loss weight, sets the
sign. The same term repairs occupancy gating (count accuracy
$0.65 \to 0.85$, hallucination rate $0.15 \to 0.05$) and, deprived of
the waveform anchor, collapses into reward hacking (SI-SDRi
$-9.4$\,dB), delimiting where such proxies can be trusted.
\item \textbf{A controlled attribution experiment and a concrete
prescription.} A joint soft-information head that bypasses explicit
separation fails at chance; controlled probes (same head, same data,
oracle vs.\ blind synchronisation as the only variable) attribute the
failure to burst-level blind carrier/phase recovery, and at $K{=}2$
even a frequency oracle fails under phase-entangled interference. The
practical conclusion for the co-frequency and PCMA communities is
direct: closing the waveform-to-bit gap requires explicit
synchronisation, equalisation, and joint-detection structure---not a
re-weighted loss.
\end{enumerate}

\paragraph{Statistical standard.}
All headline numbers are multi-seed (five seeds). Contrasts use an
exact sign-flip permutation test at the seed level and a Monte-Carlo
permutation test at the cell level, with Bonferroni control over eight
pre-declared contrasts; the $n{=}5$ power ceiling ($p \geq 0.0625$
attainable) is stated explicitly rather than smoothed over, and a
selection-bias artefact in a baseline architecture is identified and
corrected with a matched-pair control.

\paragraph{Reproducibility.}
All training, evaluation, analysis, and figure/table-generation code,
including the synthetic data generator and the deterministic test grid
(test seed 99999), is publicly available at
\url{https://github.com/BinNong/single_channel_source_seperation/tree/main/paper5_task_oriented};
no external datasets are required. The exact command behind every
number is recorded in the repository's experiment log.

\paragraph{Related submissions.}
Two related companion manuscripts are under review elsewhere and share no
results with the present paper: a closed-set separator design (the
complex CNN with squeeze-and-excitation used as a baseline family here)
at \emph{Wireless Personal Communications} (Springer), and the
open-world variable-$K$ study that this paper's benchmark and slot
separator build upon at \emph{China Communications}. The present
manuscript is fully self-contained; the cited companion manuscript
(currently listed as ``under review'') is available to the editor and
reviewers on request.

\paragraph{Originality and ethics.}
\begin{itemize}
\item The manuscript is original, has not been published elsewhere, and
is not under consideration by any other journal.
\item All authors have approved the submission; the authors have no
conflicts of interest to declare.
\item No human subjects or animal experiments are involved; all
experiments use synthetic signals from our released generator.
\end{itemize}

\paragraph{Author contributions.}
B.~Nong conceived the study and the task-oriented formulation, ran the
multi-seed experiments and statistical analyses, and drafted the
manuscript. W.~Fu contributed to the conceptualisation of the
evaluation framework. Z.~Jiang contributed resources and validation.
All authors reviewed and approved the final manuscript.

\paragraph{Funding.}
This work received no specific grant from any funding agency in the
public, commercial, or not-for-profit sectors.

\paragraph{Suggested reviewers.}
We have no specific reviewer suggestions; we trust the editorial board
to identify expert reviewers in co-frequency interference management and
machine-learning-based physical-layer signal processing.

\vspace{1em}

\noindent Thank you for considering our manuscript. We look forward to
hearing from you.

\vspace{2em}

\noindent Sincerely,\\[1.4em]
\textbf{Bin Nong}\\
Tianfu College of Southwestern University of Finance and Economics, Mianyang, China\\
\texttt{nongbin@tfswufe.edu.cn}\\[0.6em]
\textbf{Weihong Fu}\\
School of Telecommunications Engineering, Xidian University, Xi'an, Shaanxi, China\\[0.6em]
\textbf{Zhuoyun Jiang}\\
School of Health Science and Engineering, University of Shanghai for
Science and Technology, Shanghai, China

\end{document}
